% Issue 12 draft, 24 November 2026. All three Introduction paragraphs approved.
% Recruitment paragraph uses the previously agreed CIES wording.
% Final question/design/contribution paragraph: approved by the user.
% Total Introduction allowance: 0.40 pages; compiled fit remains for Issues 13--14.
\section*{Introduction}

Previous research has documented age differences in football recruitment. A 2022 CIES report found that 10.6\% of externally recruited players in Premier League squads had joined their club at age 30 or older \cite{poli2022recruitment}. In this study, we compare Premier League joining rates among eligible recorded transfers involving players aged 30+ and those under 30.

Previous research has applied sentiment analysis to football coverage from multiple newspaper sources, including the Guardian \cite{dotgrau2019representation}. Monsees (2024) found that German and Swedish newspapers constructed football talent around similar themes, including age, but gave these themes different meanings across cultural contexts \cite{monsees2024talent}. These findings motivate interpreting football coverage in its editorial, linguistic and cultural context.

We ask: \textit{How does entry into the Premier League differ by player age, and how does the tone of Guardian coverage before and after Premier League signings differ by age?} We describe eligible transfers as a network and compare Premier League joining rates for players aged 30+ and under 30. In a separate sample of selected Premier League signings, we compare Guardian sentence sentiment before and after each move to ask whether coverage tone changes differently for older and younger players. To assess the performance of the sentiment models, we compared their scores with manual labels for the same Guardian sentences. One author judged each sentence in the comparison as positive, negative or neutral, providing a human reference for the evaluation. The selected sentences from these articles form the corpus for the sentiment analysis. To describe its vocabulary, we tokenize the sentences and examine which tokens occur most often and how their frequencies vary by rank.
